% Options for packages loaded elsewhere
\PassOptionsToPackage{unicode}{hyperref}
\PassOptionsToPackage{hyphens}{url}
\documentclass[
]{article}
\usepackage{xcolor}
\usepackage[margin=1in]{geometry}
\usepackage{amsmath,amssymb}
\setcounter{secnumdepth}{-\maxdimen} % remove section numbering
\usepackage{iftex}
\ifPDFTeX
  \usepackage[T1]{fontenc}
  \usepackage[utf8]{inputenc}
  \usepackage{textcomp} % provide euro and other symbols
\else % if luatex or xetex
  \usepackage{unicode-math} % this also loads fontspec
  \defaultfontfeatures{Scale=MatchLowercase}
  \defaultfontfeatures[\rmfamily]{Ligatures=TeX,Scale=1}
\fi
\usepackage{lmodern}
\ifPDFTeX\else
  % xetex/luatex font selection
\fi
% Use upquote if available, for straight quotes in verbatim environments
\IfFileExists{upquote.sty}{\usepackage{upquote}}{}
\IfFileExists{microtype.sty}{% use microtype if available
  \usepackage[]{microtype}
  \UseMicrotypeSet[protrusion]{basicmath} % disable protrusion for tt fonts
}{}
\makeatletter
\@ifundefined{KOMAClassName}{% if non-KOMA class
  \IfFileExists{parskip.sty}{%
    \usepackage{parskip}
  }{% else
    \setlength{\parindent}{0pt}
    \setlength{\parskip}{6pt plus 2pt minus 1pt}}
}{% if KOMA class
  \KOMAoptions{parskip=half}}
\makeatother
\usepackage{color}
\usepackage{fancyvrb}
\newcommand{\VerbBar}{|}
\newcommand{\VERB}{\Verb[commandchars=\\\{\}]}
\DefineVerbatimEnvironment{Highlighting}{Verbatim}{commandchars=\\\{\}}
% Add ',fontsize=\small' for more characters per line
\usepackage{framed}
\definecolor{shadecolor}{RGB}{248,248,248}
\newenvironment{Shaded}{\begin{snugshade}}{\end{snugshade}}
\newcommand{\AlertTok}[1]{\textcolor[rgb]{0.94,0.16,0.16}{#1}}
\newcommand{\AnnotationTok}[1]{\textcolor[rgb]{0.56,0.35,0.01}{\textbf{\textit{#1}}}}
\newcommand{\AttributeTok}[1]{\textcolor[rgb]{0.13,0.29,0.53}{#1}}
\newcommand{\BaseNTok}[1]{\textcolor[rgb]{0.00,0.00,0.81}{#1}}
\newcommand{\BuiltInTok}[1]{#1}
\newcommand{\CharTok}[1]{\textcolor[rgb]{0.31,0.60,0.02}{#1}}
\newcommand{\CommentTok}[1]{\textcolor[rgb]{0.56,0.35,0.01}{\textit{#1}}}
\newcommand{\CommentVarTok}[1]{\textcolor[rgb]{0.56,0.35,0.01}{\textbf{\textit{#1}}}}
\newcommand{\ConstantTok}[1]{\textcolor[rgb]{0.56,0.35,0.01}{#1}}
\newcommand{\ControlFlowTok}[1]{\textcolor[rgb]{0.13,0.29,0.53}{\textbf{#1}}}
\newcommand{\DataTypeTok}[1]{\textcolor[rgb]{0.13,0.29,0.53}{#1}}
\newcommand{\DecValTok}[1]{\textcolor[rgb]{0.00,0.00,0.81}{#1}}
\newcommand{\DocumentationTok}[1]{\textcolor[rgb]{0.56,0.35,0.01}{\textbf{\textit{#1}}}}
\newcommand{\ErrorTok}[1]{\textcolor[rgb]{0.64,0.00,0.00}{\textbf{#1}}}
\newcommand{\ExtensionTok}[1]{#1}
\newcommand{\FloatTok}[1]{\textcolor[rgb]{0.00,0.00,0.81}{#1}}
\newcommand{\FunctionTok}[1]{\textcolor[rgb]{0.13,0.29,0.53}{\textbf{#1}}}
\newcommand{\ImportTok}[1]{#1}
\newcommand{\InformationTok}[1]{\textcolor[rgb]{0.56,0.35,0.01}{\textbf{\textit{#1}}}}
\newcommand{\KeywordTok}[1]{\textcolor[rgb]{0.13,0.29,0.53}{\textbf{#1}}}
\newcommand{\NormalTok}[1]{#1}
\newcommand{\OperatorTok}[1]{\textcolor[rgb]{0.81,0.36,0.00}{\textbf{#1}}}
\newcommand{\OtherTok}[1]{\textcolor[rgb]{0.56,0.35,0.01}{#1}}
\newcommand{\PreprocessorTok}[1]{\textcolor[rgb]{0.56,0.35,0.01}{\textit{#1}}}
\newcommand{\RegionMarkerTok}[1]{#1}
\newcommand{\SpecialCharTok}[1]{\textcolor[rgb]{0.81,0.36,0.00}{\textbf{#1}}}
\newcommand{\SpecialStringTok}[1]{\textcolor[rgb]{0.31,0.60,0.02}{#1}}
\newcommand{\StringTok}[1]{\textcolor[rgb]{0.31,0.60,0.02}{#1}}
\newcommand{\VariableTok}[1]{\textcolor[rgb]{0.00,0.00,0.00}{#1}}
\newcommand{\VerbatimStringTok}[1]{\textcolor[rgb]{0.31,0.60,0.02}{#1}}
\newcommand{\WarningTok}[1]{\textcolor[rgb]{0.56,0.35,0.01}{\textbf{\textit{#1}}}}
\usepackage{graphicx}
\makeatletter
\newsavebox\pandoc@box
\newcommand*\pandocbounded[1]{% scales image to fit in text height/width
  \sbox\pandoc@box{#1}%
  \Gscale@div\@tempa{\textheight}{\dimexpr\ht\pandoc@box+\dp\pandoc@box\relax}%
  \Gscale@div\@tempb{\linewidth}{\wd\pandoc@box}%
  \ifdim\@tempb\p@<\@tempa\p@\let\@tempa\@tempb\fi% select the smaller of both
  \ifdim\@tempa\p@<\p@\scalebox{\@tempa}{\usebox\pandoc@box}%
  \else\usebox{\pandoc@box}%
  \fi%
}
% Set default figure placement to htbp
\def\fps@figure{htbp}
\makeatother
\setlength{\emergencystretch}{3em} % prevent overfull lines
\providecommand{\tightlist}{%
  \setlength{\itemsep}{0pt}\setlength{\parskip}{0pt}}
\usepackage{bookmark}
\IfFileExists{xurl.sty}{\usepackage{xurl}}{} % add URL line breaks if available
\urlstyle{same}
\hypersetup{
  pdftitle={Stan: Introduction},
  hidelinks,
  pdfcreator={LaTeX via pandoc}}

\title{Stan: Introduction}
\author{}
\date{\vspace{-2.5em}2026-04-17}

\begin{document}
\maketitle

\subsection{Introduction}\label{introduction}

Stan is a probabilistic programming language with an R-like syntax that
compiles to highly optimised C++. Writing a model directly in Stan gives
you full control over the likelihood, the prior, and any
reparameterisation that improves sampling --- control that the
higher-level wrappers \texttt{brms} and \texttt{rstanarm} deliberately
abstract away. The cost is verbosity, but the payoff is a transparent
specification: every line in the \texttt{model} block contributes a term
to the log-posterior, and a Stan program can be archived as the
canonical statement of an analysis.

\subsection{Prerequisites}\label{prerequisites}

A working understanding of Bayesian inference, the difference between a
prior and a likelihood, and the basics of Hamiltonian Monte Carlo. You
also need a working C++ toolchain (Rtools on Windows, build-essential on
Linux) so that \texttt{rstan} or \texttt{cmdstanr} can compile the
model.

\subsection{Theory}\label{theory}

A Stan program is organised into blocks evaluated in a specific order:

\begin{itemize}
\tightlist
\item
  \texttt{data}: variables that come from the user, dimensioned and
  constrained.
\item
  \texttt{transformed\ data} (optional): deterministic transformations
  of the data computed once at compile time.
\item
  \texttt{parameters}: continuous unknowns to be inferred. Constraints
  (\texttt{\textless{}lower=0\textgreater{}}) are reparameterised
  internally, which is why discrete parameters have to be marginalised
  out.
\item
  \texttt{transformed\ parameters} (optional): deterministic functions
  of parameters, recomputed at every leapfrog step.
\item
  \texttt{model}: contributions to the log-posterior --- both priors on
  parameters and the likelihood given the data.
\item
  \texttt{generated\ quantities} (optional): posterior predictive draws,
  log-likelihoods for LOO-CV, or any deterministic quantity to monitor.
\end{itemize}

The default sampler is the No-U-Turn variant of HMC, which adapts step
size and mass matrix during a warm-up phase before producing draws.

\subsection{Assumptions}\label{assumptions}

Continuous parameters or correctly marginalised discrete latent states.
The posterior must be log-concave only locally, but pathological
geometry (funnels, ridges, multimodality) often shows up as divergent
transitions and demands reparameterisation.

\subsection{R Implementation}\label{r-implementation}

\begin{Shaded}
\begin{Highlighting}[]
\SpecialCharTok{/}\ErrorTok{/}\NormalTok{ Stan model}\SpecialCharTok{:}\NormalTok{ linear regression}
\NormalTok{data \{}
\NormalTok{  int}\SpecialCharTok{\textless{}}\NormalTok{lower}\OtherTok{=}\DecValTok{1}\SpecialCharTok{\textgreater{}}\NormalTok{ N;}
\NormalTok{  vector[N] x;}
\NormalTok{  vector[N] y;}
\NormalTok{\}}
\NormalTok{parameters \{}
\NormalTok{  real alpha;}
\NormalTok{  real beta;}
\NormalTok{  real}\SpecialCharTok{\textless{}}\NormalTok{lower}\OtherTok{=}\DecValTok{0}\SpecialCharTok{\textgreater{}}\NormalTok{ sigma;}
\NormalTok{\}}
\NormalTok{model \{}
\NormalTok{  alpha }\SpecialCharTok{\textasciitilde{}} \FunctionTok{normal}\NormalTok{(}\DecValTok{0}\NormalTok{, }\DecValTok{10}\NormalTok{);}
\NormalTok{  beta  }\SpecialCharTok{\textasciitilde{}} \FunctionTok{normal}\NormalTok{(}\DecValTok{0}\NormalTok{, }\FloatTok{2.5}\NormalTok{);}
\NormalTok{  sigma }\SpecialCharTok{\textasciitilde{}} \FunctionTok{exponential}\NormalTok{(}\DecValTok{1}\NormalTok{);}
\NormalTok{  y }\SpecialCharTok{\textasciitilde{}} \FunctionTok{normal}\NormalTok{(alpha }\SpecialCharTok{+}\NormalTok{ beta }\SpecialCharTok{*}\NormalTok{ x, sigma);}
\NormalTok{\}}
\end{Highlighting}
\end{Shaded}

\begin{Shaded}
\begin{Highlighting}[]
\FunctionTok{library}\NormalTok{(rstan)}

\FunctionTok{set.seed}\NormalTok{(}\DecValTok{2026}\NormalTok{)}
\NormalTok{d }\OtherTok{\textless{}{-}} \FunctionTok{list}\NormalTok{(}\AttributeTok{N =} \DecValTok{100}\NormalTok{, }\AttributeTok{x =} \FunctionTok{rnorm}\NormalTok{(}\DecValTok{100}\NormalTok{), }\AttributeTok{y =} \FunctionTok{rnorm}\NormalTok{(}\DecValTok{100}\NormalTok{, }\DecValTok{2} \SpecialCharTok{+} \FloatTok{1.5} \SpecialCharTok{*} \FunctionTok{rnorm}\NormalTok{(}\DecValTok{100}\NormalTok{), }\DecValTok{1}\NormalTok{))}

\CommentTok{\# Assume the Stan code is saved to "model.stan"}
\CommentTok{\# fit \textless{}{-} stan("model.stan", data = d, chains = 4, iter = 2000)}
\CommentTok{\# print(fit, pars = c("alpha", "beta", "sigma"))}
\end{Highlighting}
\end{Shaded}

\subsection{Output \& Results}\label{output-results}

\texttt{print(fit)} returns posterior mean, standard error of the mean,
posterior standard deviation, 2.5 / 25 / 50 / 75 / 97.5 \% quantiles,
effective sample size, and \(\hat R\) for every monitored parameter.
Diagnostics also include divergent-transition counts, maximum tree depth
saturation, and energy diagnostics --- any of which, if non-zero,
signals geometry trouble.

\subsection{Interpretation}\label{interpretation}

Compared with \texttt{lm()} or \texttt{glmer()} style output, Stan
output is more verbose but more explicit: there is no implicit prior, no
black-box optimisation, and the credible interval has a direct
probability interpretation conditional on the model. A reporting
sentence: ``A Stan-compiled linear regression with weakly informative
priors estimated the slope at 1.49 (95 \% CrI 1.30 to 1.69), with
\(\hat R = 1.00\) and \(n_{\text{eff}} > 3{,}500\) for every
parameter.''

\subsection{Practical Tips}\label{practical-tips}

\begin{itemize}
\tightlist
\item
  Use \texttt{brms::stancode()} or \texttt{rstanarm::stan\_glm} source
  as templates --- they are well-engineered models you can copy and
  modify.
\item
  Reparameterise hierarchical models in non-centred form
  (\(\theta_j = \mu + \tau \cdot \tilde\theta_j\) with
  \(\tilde\theta_j \sim \mathcal N(0,1)\)) to eliminate funnel-shaped
  divergences.
\item
  Prefer \texttt{cmdstanr} over \texttt{rstan} for new projects: it
  tracks upstream Stan releases, has a cleaner API, and avoids the
  binary-compatibility surprises that plague R-package builds.
\item
  Vectorise statements
  (\texttt{y\ \textasciitilde{}\ normal(mu,\ sigma)} over \texttt{for}
  loops); the compiler optimises vectorised forms more aggressively.
\item
  Place posterior predictive draws in \texttt{generated\ quantities},
  not the \texttt{model} block, so they do not enter the log-density
  evaluation and slow sampling.
\end{itemize}

\subsection{R Packages Used}\label{r-packages-used}

\texttt{rstan} or \texttt{cmdstanr} to compile and sample from Stan
programs, \texttt{posterior} for working with draw arrays,
\texttt{bayesplot} for diagnostic and posterior plots, and \texttt{loo}
for leave-one-out cross-validation.

\subsection{For Reviewers}\label{for-reviewers}

\textbf{What to look for in a paper using this method.}

\begin{itemize}
\tightlist
\item
  \textbf{Common misapplications.}
\item
  \textbf{Diagnostics that should be reported but often aren't.}
\item
  \textbf{Red flags in tables and figures.}
\item
  \textbf{What to verify.}
\item
  \textbf{What an adequate Methods paragraph must contain.}
\end{itemize}

\subsection{See also --- labs in this
chapter}\label{see-also-labs-in-this-chapter}

\begin{itemize}
\tightlist
\item
  \href{labs/lab-bayesian-thinking-control-vs-decision-error.Rmd}{Bayesian
  thinking; α control vs decision error}
\item
  \href{labs/lab-brms-stan-loo-hierarchical-models.Rmd}{brms / Stan,
  LOO, hierarchical models}
\end{itemize}

\end{document}
